\section{Telemetry, Runtime Tuning, and Scenarios}

\subsection{Telemetry}

The core emits numeric telemetry tuples through the platform vtable. Platforms decide whether those tuples are sent over UDP, USB-CDC, UART, CAN, or written as CSV.

\begin{lstlisting}[caption={Telemetry frame sketch}]
u32 timestamp_ms
u16 signal_id
u16 reserved
i32 value
\end{lstlisting}

\subsection{Runtime Tuning}

Runtime tuning is a bench and development feature. Commands may adjust PID gains, setpoints, trip points, curve points, and fault injection state, but all values remain bounded by compile-time or configuration limits.

\subsection{Scenario Runner}

Scenario files should make benchmark runs reproducible. For speed-aware scenarios, the runner may command \texttt{car-can-emulator} over TCP while \texttt{thermal-core} observes only CAN traffic.

\subsection{To Fill In Later}

\begin{itemize}[leftmargin=*]
    \item final command frame format;
    \item \texttt{thermalcore-probe} CLI examples;
    \item \texttt{thermalcore-tune} CLI examples;
    \item first three canonical scenario files.
\end{itemize}

\endinput
